\section{物理层}

\subsection{物理层的基本概念}

\begin{mydef}{物理层}{physical-layer}
物理层是 OSI 的最底层，为数据链路层提供透明的比特流传输。它规定：
\begin{itemize}
    \item \textbf{机械特性}：连接器的形状、尺寸、引脚数和排列；
    \item \textbf{电气特性}：信号电平范围、传输速率、传输距离；
    \item \textbf{功能特性}：每个引脚信号的物理意义；
    \item \textbf{规程特性}：信号的工作时序。
\end{itemize}
物理层传输的是比特流，不关心比特的含义（帧结构、地址等由上层处理）。
\end{mydef}

\begin{attention}
\textbf{408 辨析：} 物理层传输的是「比特」（bit），数据链路层传输的是「帧」（Frame）。集线器（Hub）工作在物理层，只转发比特不识别帧结构。交换机（Switch）工作在数据链路层，能识别 MAC 地址并转发帧。
\end{attention}

\subsection{数据通信基础}

\begin{conclusion}{数据通信系统模型}{comm-model}
\[
\text{信源} \to \text{发送器} \to \boxed{\text{信道}} \to \text{接收器} \to \text{信宿}
\]
\vspace{-8pt}
\begin{itemize}
    \item \textbf{信源/信宿}：产生/接收数据的设备。
    \item \textbf{发送/接收器}：编码/解码（调制解调器 Modem 即为典型的发送+接收器）。
    \item \textbf{信道}：信号的传输介质。分为有线（双绞线、同轴电缆、光纤）和无线（电磁波）。
\end{itemize}
\end{conclusion}

\begin{formula}{通信的基本术语}{comm-terms}
\begin{itemize}
    \item \textbf{数据 vs 信号}：数据是传送信息的实体，信号是数据的电气/电磁表现。
    \item \textbf{模拟信号 vs 数字信号}：连续变化 vs 离散取值。
    \item \textbf{基带传输}：直接传输数字信号（不经过调制）。
    \item \textbf{频带传输}：将数字信号调制为模拟信号后传输。
    \item \textbf{单工/半双工/全双工}：单向/双向交替/双向同时。
    \item \textbf{码元（Symbol）}：用一个固定时长的信号波形表示一个 $k$ 进制数字，代表 $k$ 种不同离散值。
\end{itemize}
\end{formula}

\subsection{奈奎斯特（Nyquist）准则与香农（Shannon）定理}

\begin{conclusion}{奈奎斯特准则（无噪声信道）}{nyquist}
在理想低通（无噪声）信道中，极限码元传输速率：
\[
B_{\text{max}} = 2W\;\text{（Baud）}
\]
其中 $W$ 为信道带宽（Hz）。若每个码元携带 $\log_2 V$ 比特（$V$ 为离散电平数），则极限数据率：
\[
C_{\text{max}} = 2W \log_2 V\;\text{（bps）}
\]
\textbf{奈奎斯特准则给出的是无噪声下的符号率上限。}
\end{conclusion}

\begin{conclusion}{香农定理（有噪声信道）}{shannon}
带宽受限且有高斯白噪声的信道中，极限数据率：
\[
C = W \log_2\!\left(1 + \frac{S}{N}\right)\;\text{（bps）}
\]
其中 $S/N$ 为信噪比（信号功率与噪声功率之比）。通常使用分贝（dB）表示：
\[
\text{信噪比(dB)} = 10\log_{10}\!\left(\frac{S}{N}\right)
\]

\textbf{香农定理给出的是信息传输速率的上限，无论使用何种编码方式都无法超越。}
\end{conclusion}

\begin{attention}
\textbf{408 综合计算必考：} 给定带宽和信噪比，先由奈奎斯特准则算符号率上限，再由香农定理算信息速率上限，取两者的\textbf{最小值}作为实际极限速率。给定 dB 值时需先换算为 $S/N$ 比值。
\end{attention}

\subsection{编码与调制}

\begin{conclusion}{数字数据的数字信号编码}{digital-encoding}
\begin{tablebox}{常见数字编码}
\centering
\begin{tabular}{|l|l|}
\hline
编码方式 & 规则 \\
\hline
不归零编码（NRZ）& 高电平=1，低电平=0；无同步信号 \\
归零编码（RZ）& 每个比特中间跳变回零 \\
曼彻斯特编码 & 比特中间跳变：高$\to$低=1，低$\to$高=0；自同步 \\
差分曼彻斯特 & 比特开始有无跳变：无跳变=1，有跳变=0；中间总有跳变 \\
\hline
\end{tabular}
\end{tablebox}
\textbf{曼彻斯特编码：}每个比特的中间都有跳变（既作时钟又作数据），波特率是数据率的 2 倍。802.3 以太网使用曼彻斯特编码。

\textbf{差分曼彻斯特：}抗干扰性更强，常用于令牌环网。
\end{conclusion}

\begin{conclusion}{数字调制技术}{digital-modulation}
将数字信号调制到模拟载波上传输：
\begin{itemize}
    \item \textbf{调幅（ASK）}：改变振幅表示 0/1。
    \item \textbf{调频（FSK）}：改变频率。
    \item \textbf{调相（PSK）}：改变相位。QPSK 用 4 个相位，每码元携带 2 bit。
    \item \textbf{QAM（正交振幅调制）}：同时改变振幅和相位，每码元携带更多比特。
\end{itemize}
\end{conclusion}

\subsection{传输介质}

\begin{conclusion}{常见传输介质}{transmission-media}
\begin{tablebox}{有线传输介质对比}
\centering
\begin{tabular}{|l|c|l|}
\hline
介质 & 典型带宽/距离 & 特点 \\
\hline
双绞线 & 100m (100Mbps-10Gbps) & 便宜，易受干扰；Cat5e/Cat6 \\
同轴电缆 & 更长距离 & 抗干扰优于双绞线 \\
光纤（单模）& 数十公里以上 & 带宽极大，损耗极低，最昂贵 \\
光纤（多模）& 几百米-数千米 & 带宽大，成本低于单模 \\
\hline
\end{tabular}
\end{tablebox}

\textbf{408 常考：}光纤利用光的\textbf{全反射}原理传输，单模光纤芯径小（~9$\mu$m）传输一个模式，多模光纤芯径大（~50$\mu$m）传输多个模式，存在模间色散。
\end{conclusion}

\subsection{物理层设备}

\begin{conclusion}{中继器与集线器}{repeater-hub}
\begin{itemize}
    \item \textbf{中继器（Repeater）}：对衰减的信号进行放大和整形，延长网络覆盖距离。仅有两个端口，工作在物理层。
    \item \textbf{集线器（Hub）}：多端口中继器。从一个端口收到的信号放大整形后广播到所有其他端口。所有端口共享同一冲突域。
\end{itemize}

\textbf{5-4-3 规则（以太网）：}网络中最多 5 个网段，其中最多 4 个中继器/集线器，最多 3 个网段可以连接计算机。
\end{conclusion}

\subsection{408 高频考点速查}

\begin{conclusion}{物理层 — 408常考点}{cn2-keypoints}
\begin{enumerate}
    \item \textbf{奈奎斯特 + 香农综合计算}：给定带宽、电平数 $V$、信噪比 dB，求最大数据率。
    \item \textbf{曼彻斯特编码 vs 差分曼彻斯特}：波形图的手工绘制/识别，比特率和波特率的关系。
    \item \textbf{物理层设备}：中继器/集线器不隔离冲突域。
    \item \textbf{码元与比特的关系}：一个码元携带 $\log_2 V$ 比特。
    \item \textbf{电路交换/报文交换/分组交换的时延计算}（与第 1 章交叉）。
\end{enumerate}
\end{conclusion}

\newpage
\section{习题}

\subsection{基础过关题}

\begin{enumerate}

\item \textbf{（单项选择题）} 物理层传输的数据单元是
\begin{center}
\begin{tabular}{ll}
(A) 帧 & (B) 分组 \\[6pt]
(C) 比特 & (D) 报文
\end{tabular}
\end{center}

\item \textbf{（单项选择题）} 曼彻斯特编码的波特率与比特率的关系是
\begin{center}
\begin{tabular}{ll}
(A) 波特率 = 比特率 & (B) 波特率 = 2$\times$比特率 \\[6pt]
(C) 比特率 = 2$\times$波特率 & (D) 两者无关
\end{tabular}
\end{center}

\item \textbf{（填空题）} 若某信道带宽为 3 kHz，信噪比为 30 dB，则根据香农定理，该信道的极限数据率为 \underline{\qquad} bps（保留整数）。

\item \textbf{（简答题）} 简述集线器（Hub）和交换机（Switch）在 OSI 参考模型中所处的层次以及主要功能差异。

\item \textbf{（计算题）} 若某信道带宽为 4 kHz，采用 16 种不同相位的 QPSK 调制技术（即每个码元有 16 种不同的离散值）。在无噪声条件下，该信道的最大数据率是多少？

\end{enumerate}

\subsection{真题风格与综合训练}

\begin{enumerate}[resume]

\item \textbf{（真题风格，单项选择题）} 以下关于曼彻斯特编码和差分曼彻斯特编码的说法，正确的是
\begin{center}
\begin{tabular}{ll}
(A) 两者均无自同步能力 \\[4pt]
(B) 差分曼彻斯特编码中，比特中间的跳变仅用于同步 \\[4pt]
(C) 曼彻斯特编码的波特率等于比特率 \\[4pt]
(D) 差分曼彻斯特编码的抗干扰性更好
\end{tabular}
\end{center}

\item \textbf{（真题风格，综合题）} 某信道的带宽为 3 kHz，采用四相调制（4 种相位，每种相位一个振幅），信噪比为 20 dB。
\begin{enumerate}
    \item[(1)] 根据奈奎斯特准则，该信道的最大数据率是多少？
    \item[(2)] 根据香农定理，该信道的极限数据率是多少？
    \item[(3)] 该信道实际可达到的最大数据率是多少？说明理由。
\end{enumerate}

\item \textbf{（真题风格，综合题）} 主机 A 和主机 B 之间通过一个集线器和两段 100 m 的双绞线连接，端到端传输速率为 100 Mbps，电磁波在双绞线上的传播速率为 $2\times10^8$ m/s。A 向 B 发送一个长度为 1000 字节的帧。
\begin{enumerate}
    \item[(1)] 假设采用电路交换方式（忽略电路建立时间），求 A 开始发送到 B 完全收完的总时延；
    \item[(2)] 若采用分组交换，A 先将帧分为两个 500 字节的分组依次发送。集线器采用存储转发方式，求总时延；
    \item[(3)] 对比两种方式的结果。
\end{enumerate}

\end{enumerate}

\vspace{8pt}
\noindent\textbf{拓展阅读：}
\begin{itemize}
    \item Kurose \& Ross《计算机网络：自顶向下方法》第 1 章（1.4-1.5 节）——时延、丢包和吞吐量的定量分析，含奈奎斯特/香农的直观解释；
    \item 谢希仁《计算机网络》第 2 章——物理层传输介质和编码技术的标准国内表述。
\end{itemize}
